%!TEX root=../assign

\subsection{Problem 1}
\begin{enumerate}[label=(\Roman*)]
    \item \textbf{Prove that $\sqrt{3}$ is irrational.}
    
    We need to prove a lemma first before we prove the proposition. 
     \begin{lemma} \label{lemma 1}
            For all $n \in \Z$, if $3 \mid n^2$, then $3 \mid n$. 

            \begin{proof}
                Suppose $3 \nmid n$. There are two cases to discuss.  
                
                case 1: $n = 3k + 1$

                According to the definition of divisibility, we have 
                \begin{equation*}
                    n = 3k + 1 \quad \exists k \in \Z. 
                \end{equation*} 

                Taking the square of both sides, it follows that 
                \begin{equation*}
                    n^2 = \paren{3k + 1}^2 = 9k^2 + 6k + 1 =  3c + 1 \quad c = 3k^2 + 2k \in \Z.
                \end{equation*}

                thus, $3 \nmid n^2$.\\

                case 2: $n = 3k + 2$

                According to the definition of divisibility, we have 
                \begin{equation*}
                     \exists k \in \Z \text{ s.t } n = 3k + 2  
                \end{equation*} 

                Taking the square of both sides, it follows that 
                \begin{equation*}
                    n^2 = \paren{3k + 2}^2 = 9k^2 + 12k + 4 =  3c + 1 \quad c = 3k^2 + 4k + 1 \in \Z.
                \end{equation*}

                thus, $3 \nmid n^2$.
            \end{proof}
        \end{lemma}

        Here is the proof of the question:
        \begin{proof}
            Suppose $\sqrt{3} \in \Q$. 

            Then, according to the definition of rational numbers, 
            \begin{equation*}
                \exists m, n \in \Z, n \neq 0 \quad \text{ s.t } \quad \sqrt{3} = \frac{m}{n}. 
            \end{equation*}

            We assume that $m$ and $n$ have no common factor other than 1. Taking the square of the both sides, we get 
            \begin{equation*}
                3 = \frac{m^2}{n^2} \quad \implies \quad m^2 = 3n^2, 
            \end{equation*}

            which implies that $3 \mid m$ according to the lemma \ref{lemma 1}. 
            According to the definition of divisibility, 
            \begin{equation*}
                \exists k \in \Z \quad \text{ s.t } \quad  m = 3k. 
            \end{equation*}

            Subtituting $m = 3k$ back into $m^2 = 3n^2$, we have 
           \begin{align*}
            n^2 &= (3k)^2 / 3 \\
            &= 9k^2 / 3\\
            &= 3d, d = k^2 \in \Z. 
           \end{align*}

           It means that $3 \mid n$. 

           Since $3 \mid n$ and $3 \mid m$, it contradicts the assumption that $n$ and $m$ have no common factor other than 1. 
           
           Therefore, $\sqrt{3} \in \R \setminus\Q$. 

        \end{proof}

    \item \begin{proof}
        Suppose $a, b \in \R$ with $a > 0$ and $b > 0$.

        According to the given statement, we have
        \begin{equation*}
            \paren{\forall c > 0}\paren{\exists x > 0} x^2 = c.
        \end{equation*}

        It follows that there exist positive real numbers $a'$ and $b'$ such that
        \begin{equation*}
            \paren{a'}^2 = a
            \quad \text{and} \quad
            \paren{b'}^2 = b.
        \end{equation*}

        By the ordered field properties, we have
        \begin{equation*}
            \paren{a' - b'}^2 \geq 0.
        \end{equation*}

        for both $a' - b' \geq 0$ and $b'-a' \geq 0$. 

        Expanding the left side, we get
        \begin{equation}
            \paren{a'}^2 - 2a'b' + \paren{b'}^2 \geq 0
            \quad \implies \quad
            \frac{\paren{a'}^2 + \paren{b'}^2}{2} \geq a'b'.
            \label{eq 1.2.1}
        \end{equation}

        By substituting $\paren{a'}^2 = a$ and $\paren{b'}^2 = b$ into equation $\eqref{eq 1.2.1}$, we have
        \begin{equation*}
            \frac{a + b}{2} \geq a'b'.
        \end{equation*}

        Since $a' > 0$ and $b' > 0$, it follows from the ordered field properties that
        \begin{equation*}
            a'b' > 0 \cdot b' = 0.
        \end{equation*}

        Moreover,
        \begin{equation*}
            \paren{a'b'}^2
            = \paren{a'}^2 \paren{b'}^2
            = ab > 0 \cdot a'b' = 0.
        \end{equation*}

        Since $a'b' > 0$ and $ab > 0$, therefore,
        \begin{equation*}
            a'b' = \sqrt{ab}.
        \end{equation*}

        As a result,
        \begin{equation*}
            \frac{a + b}{2} \geq \sqrt{ab}.
        \end{equation*}

        Thus, the inequality has been proved.
    \end{proof}

    \textbf{Reference:} The ordered field properties were consulted in V. A. Zorich, \textit{Mathematical Analysis I}, Section 2.2.
    Prove style is modified based on the suggestions from my dear father, Changwei Hu. 


\end{enumerate}